\documentclass[12pt]{article}

\usepackage{amsmath}
\usepackage{amssymb}
\usepackage{graphicx}
\usepackage{hyperref}
\usepackage[margin=1in]{geometry}
\usepackage{natbib}
\usepackage{xcolor}
\usepackage{tcolorbox}
\usepackage{booktabs}
\usepackage{tabularx}
\usepackage{lineno}
\linenumbers
\tcbuselibrary{breakable,skins}
\bibliographystyle{aasjournal}

\newcommand{\HD}{\mathrm{H_D}}
\newcommand{\HDp}{\mathrm{H_D^+}}
\newcommand{\HDm}{\mathrm{H_D^-}}
\newcommand{\HDtwo}{\mathrm{H_{D,2}}}
\newcommand{\HDtwop}{\mathrm{H_{D,2}^+}}
\newcommand{\pD}{\mathrm{p_D}}
\newcommand{\eD}{\mathrm{e_D}}
\newcommand{\gD}{\gamma_D}
\newcommand{\alphaD}{\alpha_D}
\newcommand{\eV}{\,\mathrm{eV}}
\newcommand{\cmg}{\,\mathrm{cm^2\,g^{-1}}}
\newcommand{\Msun}{\,M_\odot}
\newcommand{\smx}{\sigma/m}
\newcommand{\ra}{r_\alpha}
\newcommand{\rmm}{r_m}      % r_m collides with \rm; use \rmm inside math
\newcommand{\rM}{r_M}
\newcommand{\aD}{a_{0,D}}

\title{Molecular Chemistry for Atomic Dark Matter}
\author{}
\date{\today}

\begin{document}

\maketitle
\tableofcontents
\newpage

\begin{tcolorbox}[colback=yellow!10!white, colframe=orange!75!black,
                  title={\bfseries Disclaimer}]
These notes were compiled by an AI assistant (Claude, Anthropic) and have
\textbf{not been reviewed by a human domain expert}.  They may contain
errors in derivations, physical reasoning, or numerical estimates.
Cross-check key equations against the primary literature before use.
\end{tcolorbox}
\vspace{0.5em}

%----------------------------------------------------------------------
\section{Motivation and scope}
\label{sec:motivation}

Standard cold dark matter (CDM) is treated as a single collisionless
species with no internal structure.  A minimal deviation is to allow
the dark sector to have its own gauge interaction, so that two
oppositely charged fermions can bind into ``atoms'' analogous to
hydrogen, and those atoms can dissipate energy through excitation and
radiative de-excitation.  This class of models --- \emph{atomic dark
matter} \citep{GoldbergHall1986, Ackerman2009, Kaplan2010,
CyrRacineSigurdson2013, Fan2013} --- is not proposed to solve every
cosmological problem, but it is the simplest playground in which the
consequences of a dissipative dark sector can be worked out
quantitatively.

The interest is twofold.  First, if the dark sector cools, it can
fragment and form compact objects: dark protostars, dark stars, and
dark black holes.  A subsolar-mass dark black hole binary is a
distinctive gravitational-wave signature that LIGO can constrain, with
no standard-model analogue \citep{Shandera2018, Singh2020}.  Second,
even before compact-object formation, atomic DM alters the linear
matter power spectrum through the analogue of Silk damping (dark
diffusion) and baryon acoustic oscillations (dark acoustic
oscillations, DAOs), which are directly probed by galaxy surveys
\citep{CyrRacineSigurdson2013}.

This note summarises the machinery for treating atomic-DM chemistry
developed in the trilogy of \citet{Ryan2022a}, \citet{Gurian2022}, and
\citet{Ryan2022b}, hereafter Paper~I, II, and III.

%----------------------------------------------------------------------
\section{The atomic-DM model}
\label{sec:model}

Introduce two oppositely charged fermions --- ``dark proton'' $\pD$ of
mass $M$ and ``dark electron'' $\eD$ of mass $m$ --- interacting via a
massless $U(1)$ dark photon $\gD$ with fine-structure constant
$\alphaD$.  Take the hierarchy $M \gg m$, exactly as in the standard
model (SM).  The dark sector then has:
%
\begin{itemize}
  \item A dark Bohr radius $\aD = 1/(m \alphaD)$ (setting
        $\hbar=c=1$; conversion in Appendix~\ref{app:natural_units}).
  \item A dark Rydberg (binding energy)
        $E_{H_D} = \tfrac{1}{2}m\alphaD^2$.
  \item Bound states $\HD$ (dark hydrogen), $\HDtwo$ (dark H$_2$),
        their ions $\HDp$, $\HDm$, $\HDtwop$, plus free protons and
        electrons.
\end{itemize}
%
Both $\aD$ and $E_{H_D}$ are derived from the Schr\"odinger equation
in Appendix~\ref{app:hydrogenic}: $\aD$ from dimensional analysis
(the only length that can be built from $m$, $\alphaD$, $\hbar$, $c$)
and $E_{H_D}$ from the virial theorem
$E_{\rm bound} = \tfrac{1}{2}\langle V\rangle$.  See
Appendix~\ref{app:U1} for what ``$U(1)$'' and gauge coupling mean.
%
Cosmologically, one additional parameter appears: the ratio of the
dark-photon background temperature to the CMB temperature,
%
\begin{equation}
  \xi \equiv \frac{T_{\gD}}{T_\gamma}.
\end{equation}
%
$N_{\rm eff}$ constraints from BBN and CMB require $\xi \lesssim 0.4$.
The natural default in Paper~II is $\xi = 0.37$; a colder dark sector
($\xi \sim 0.01$) is the default in Paper~III, chosen so that DAOs are
suppressed to observationally acceptable levels.

The four numbers $\{M, m, \alphaD, \xi\}$ fully specify the model; SM
hydrogen physics corresponds to $M = m_p$, $m = m_e$,
$\alphaD = \alpha_{\rm SM} = 1/137$, $\xi = 1$.

\subsection{Why the dark plasma stays quasi-neutral}
\label{sec:quasi_neutral}

A subtle point: the dark sector has a heavy species ($\pD$) and a
light species ($\eD$) with different masses.  Under gravity alone the
heavier species would sink to the halo centre while the lighter one
would not --- ``sedimentation'' would split the plasma into
$\pD$-rich core plus $\eD$-rich envelope.  Why doesn't this happen?

Because the $U(1)$ Coulomb force is much stronger than gravity.
Estimate the ratio of the two forces between one $\pD$ and one $\eD$:
%
\begin{equation}
  \frac{F_{\rm Coulomb}}{F_{\rm grav}}
    = \frac{\alpha_D/r^2}{G M m/r^2}
    = \frac{\alpha_D}{G M m}.
\end{equation}
%
For SM values ($\alpha \sim 10^{-2}$, $G m_p m_e \sim 10^{-45}$ in
natural units), the ratio is $\sim 10^{43}$: gravity is negligible
compared to the electromagnetic attraction.  For the dark sector
Ryan+2022 point out that the analogous inequality
%
\begin{equation}
  \alpha_D \gg G M m
  \label{eq:quasi_neutral}
\end{equation}
%
holds trivially throughout the parameter space they consider.  As a
result, any charge separation of order $\lambda_{\rm sep}$ builds up a
restoring electric field $E \sim g_D n_D \lambda_{\rm sep}$
(Gauss's law), which pulls the two species back together far faster
than gravity can pull them apart.

This has three consequences:
%
\begin{itemize}
  \item \textbf{Dark plasma is neutral on halo scales.}
        $n_{\pD}(\mathbf{x}) \approx n_{\eD}(\mathbf{x})$ everywhere.  The dark
        sector behaves as a single-species fluid for dynamical
        purposes, not a two-fluid $(\pD, \eD)$ system.
  \item \textbf{Two-fluid ambipolar treatment is unnecessary.}
        Roy+2024's use of a \emph{single} hydro fluid for aDM (not two
        separate fluids for dark protons and dark electrons) is
        justified by \eqref{eq:quasi_neutral}: the two species are
        pinned to the same density and velocity by their mutual
        Coulomb force.
  \item \textbf{Dark plasma oscillations at $\omega_p$.}  On scales
        much shorter than the Debye length
        $\lambda_D = \sqrt{T/(4\pi\alpha_D n_D)}$, small charge
        separations oscillate at the dark plasma frequency
        $\omega_p = \sqrt{4\pi\alpha_D n_D/m}$ instead of running
        away.  Same physics as SM baryon-photon coupling before
        recombination.
\end{itemize}

The quasi-neutrality argument is the microscopic reason that atomic
dark matter can be treated as a single dissipative fluid in
simulations --- the underlying $\pD, \eD$ 2-fluid structure is
integrated out at scales larger than the Debye length, which is
always much smaller than any halo.

%----------------------------------------------------------------------
\section{Rescaling from Standard Model hydrogen}
\label{sec:rescaling}

The central technical trick of Paper~I is that every rate coefficient
in dark hydrogen chemistry can be obtained by \emph{rescaling} the
corresponding SM rate, without redoing the quantum-mechanical
calculation.  This works because in the Born--Oppenheimer approximation
the molecular potential depends only on the electronic energy scale
$E_H$ and the Bohr radius $a_0$, whose dependences on the fundamental
parameters are known.  Define the three dimensionless ratios
%
\begin{equation}
  \ra \equiv \frac{\alphaD}{\alpha_{\rm SM}},
  \qquad
  \rmm \equiv \frac{m}{m_e},
  \qquad
  \rM \equiv \frac{M}{m_p}.
\end{equation}
%
Every length, energy, and rate then scales as some power of $(\ra,
\rmm, \rM)$ times its SM value.  For example:
%
\begin{align}
  a_{0,D} &= a_0\,\ra^{-1}\,\rmm^{-1},
  \label{eq:aD}\\
  E_{H_D} &= E_H\,\ra^2\,\rmm,
  \label{eq:EHD}\\
  E_{\rm rot,D} &= E_{\rm rot,SM}\,\ra^2\,\rmm^2\,\rM^{-1},
  \label{eq:Erot}\\
  E_{\rm vib,D} &= E_{\rm vib,SM}\,\ra^2\,\rmm^{3/2}\,\rM^{-1/2}.
  \label{eq:Evib}
\end{align}
%
Rotational levels scale like $\hbar^2/I \propto 1/M X_0^2$ (moment of
inertia of a rigid rotor); vibrational levels scale like
$\hbar\sqrt{k/M}$ with the spring constant $k \propto V_0/b^2 \propto
m^3\alpha^4$ read off the Morse potential curvature at the well
minimum.  \citet{Ryan2022a} Table~1 gives the complete list of
rescalings for reaction rates, Einstein $A$ coefficients (the
homonuclear $\HDtwo$ transitions are electric-quadrupole because the
dipole moment vanishes, exactly as in SM), and polarizabilities.

The physical statement is: as long as the dark sector is weakly coupled
($\alphaD \ll 1$) and mass-hierarchical ($m \ll M$), the entire
molecular physics is a one-parameter family in
$(\ra, \rmm, \rM)$ around the SM value.

%----------------------------------------------------------------------
\section{Dark recombination and cosmological freeze-out}
\label{sec:recomb}

Paper~II adapts the standard SM recombination code
\texttt{Recfast++} \citep{Peebles1968, Seager1999} to the dark
sector.  The Peebles effective three-level structure ($1s$, $2s$,
$2p$) is preserved, but every rate is rescaled per
Section~\ref{sec:rescaling} and the molecular processes of
\citet{GalliPalla1998} are added:
%
\begin{align}
  \pD + \eD &\rightarrow \HD + \gD
     &&\text{(radiative recombination)},\\
  \HD + \gD &\rightarrow \HDm + \gD^{\rm em}
     &&\text{(H$^-$ analog formation)},\\
  \HDm + \HD &\rightarrow \HDtwo + \eD
     &&\text{(H$^-$ channel to H$_2$)},\\
  \HD + \pD &\rightarrow \HDtwop + \gD
     &&\text{(H$_2^+$ channel)},\\
  \HDtwop + \HD &\rightarrow \HDtwo + \pD
     &&\text{(charge transfer to H$_2$)}.
\end{align}
%
Ionisation fractions $x_{\pD}$, $x_{\eD}$, $x_{\HDp}$, $x_{\HDtwo}$,
$x_{\HDtwop}$ are evolved in redshift against Hubble expansion.

\subsection{Why recombination happens at $T\ll E_H$}
\label{sec:saha_factor}

Naively you would expect a plasma to recombine when the photon
temperature $T$ drops below the atomic binding energy $E_H$: at
$T\lesssim E_H$ there are not enough energetic photons to ionise
newly-formed atoms, so recombination wins.  But actual (SM)
recombination happens at $T_\gamma \simeq 0.3\eV$, some \emph{40 times
smaller} than $E_H = 13.6\eV$.  The reason is the enormous
photon-to-baryon ratio.

\paragraph{Saha equation.}  In thermal equilibrium, the reaction
$e^- + p \rightleftharpoons H + \gamma$ satisfies detailed balance,
which for non-relativistic species gives the \emph{Saha equation}
%
\begin{equation}
  \frac{n_e n_p}{n_H}
    = \left(\frac{m_e T}{2\pi}\right)^{3/2}
      e^{-E_H/T},
  \label{eq:saha_dm}
\end{equation}
%
where $n_H$ is neutral hydrogen number density.  Divide both sides by
the total baryon density $n_b = n_p + n_H$ and use charge neutrality
$n_e = n_p = x_e n_b$:
%
\begin{equation}
  \frac{x_e^2}{1-x_e}
    = \frac{1}{n_b}\!\left(\frac{m_e T}{2\pi}\right)^{3/2}\!
      e^{-E_H/T}.
  \label{eq:saha_xe}
\end{equation}
%
Because $n_b/n_\gamma \simeq 6\times 10^{-10}$ (very few baryons per
photon), the prefactor $(m_e T)^{3/2}/n_b$ is enormous --- of order
$10^{16}$ at $T\sim E_H$.  To make $x_e$ drop to a small value the
Boltzmann suppression $e^{-E_H/T}$ must overcome this factor, which
requires
%
\begin{equation}
  \frac{E_H}{T_{\rm rec}}
    \simeq \ln\!\left[\frac{1}{n_b}\!\left(\frac{m_e T}{2\pi}\right)^{3/2}\right]
    \simeq 40
  \qquad\text{(SM)}.
\end{equation}
%
Hence ``$T_{\rm rec} = E_H / \text{Saha factor}$'' with Saha factor
$\simeq 40$: recombination happens well below the naive $T\sim E_H$
temperature because photons still outnumber matter by nine orders of
magnitude even after Boltzmann suppression.  Numerically
$T_{\rm rec}^{\rm SM}\simeq 0.34\eV$, corresponding to redshift
$z_{\rm rec}\simeq 1090$.

\paragraph{Saha fails; Peebles closes it.}  Equation \eqref{eq:saha_xe}
holds only \emph{in equilibrium}, which is maintained as long as
recombination happens faster than the Hubble expansion.  Around
$T\sim E_H/40$, the recombination rate drops below $H$ and $x_e$
freezes out at $\sim 10^{-3}$ instead of Saha's $\sim 10^{-15}$.
Solving this out-of-equilibrium freeze-out requires the Peebles
effective 3-level rate equations (Appendix~\ref{app:cosmo_primer}).
\texttt{Recfast++} \citep{Seager1999} does this fast; the atomic-DM
port in Paper~II is essentially the same code with rescaled rate
coefficients and dark photon temperature $T_{\gD} = \xi T_\gamma$.

\paragraph{Dark recombination.}  For dark hydrogen, use dark binding
energy $E_{H_D} = \ra^2 \rmm E_H$ (with $\ra = \alpha_D/\alpha_{\rm
EM}$, $\rmm = m/m_e$) and dark radiation temperature $T_{\gD} =
\xi T_\gamma$.  The Saha factor is now
%
\begin{equation}
  \frac{E_{H_D}}{T_{\rm rec}^D}
    \simeq \ln\!\left[\frac{1}{n_{b,D}}\!\left(\frac{m\,T_{\gD}}{2\pi}\right)^{3/2}\right],
\end{equation}
%
which depends on both the intrinsic dark binding energy and on the
dark radiation temperature (via $\xi$).  Freeze-out abundances
(residual free electrons that later catalyse molecule formation)
shift substantially with $(\alpha_D, m, M, \xi)$; that is why the
Paper~II numerical port matters.

%----------------------------------------------------------------------
\section{Dark diffusion and dark acoustic oscillations}
\label{sec:dao}

Before recombination, dark photons and dark electrons/protons are
tightly coupled, so the dark plasma behaves as a single fluid with
finite sound speed $c_{s,D}$.  Two effects imprint on the linear
power spectrum:
%
\begin{itemize}
  \item \textbf{Dark diffusion (Silk-analogue) damping.}
        Photons diffuse with a Compton mean free path
        $\lambda_{\gD} \propto 1/(\alphaD^2 n_{\eD})$;
        the associated diffusion scale $k_D^{\rm diff}$ sets an
        exponential suppression $\exp(-k^2/k_D^{\rm diff\,2})$ of
        power on smaller scales.
  \item \textbf{Dark acoustic oscillations (DAOs).}
        The finite sound speed produces standing-wave features
        (BAO-analogue) at the dark sound horizon at decoupling.
        On scales $k > k_{\rm DAO}$ the power spectrum shows
        oscillatory ringing that damps into the diffusion cutoff.
\end{itemize}
%
Together these effects give an atomic-DM transfer function with a
warm-dark-matter-like cutoff superposed with dark BAO.  In the
low-$\xi$ regime the DAO amplitude is small and the phenomenology
resembles WDM with an effective free-streaming scale set by the dark
diffusion length instead of thermal velocity.

The observational constraint $\xi \lesssim 0.4$ from $N_{\rm eff}$ is
weaker than the constraint on DAOs from galaxy surveys, which pushes
$\xi$ downward once atomic DM is a non-negligible fraction of the
total DM \citep{CyrRacineSigurdson2013}.

Paper~II computes the atomic-DM halo mass function at formation
redshift using the diffusion cutoff as an effective ansatz, in analogy
with the WDM mass function calibrated on $N$-body simulations.
Subsequent evolution --- crucially including radiative cooling and
fragmentation --- alters the mass function further at low redshift,
and is the target of Paper~III's zone-collapse simulations.

%----------------------------------------------------------------------
\section{Cooling and compact-object formation}
\label{sec:cooling}

Once a dark halo virialises with temperature $T_{\rm vir}$, whether it
collapses to a compact object depends on whether it can radiate away
its virial kinetic energy faster than a Hubble time.  In SM primordial
gas, the dominant coolant below $10^4\,{\rm K}$ is molecular hydrogen
line cooling.  Below $\sim 100\,{\rm K}$, cooling stalls (the H$_2$
rotational lines get too widely spaced), setting the Jeans mass and
hence the mass of the first stars.

The atomic-DM analogue of this argument, with all rates rescaled, is
Paper~I's application in Section~3 of \citet{Ryan2022a}: rovibrational
$\HDtwo$ cooling drives the gas to a temperature floor set by the
lowest available line spacing, which scales as $\ra^2 \rmm^2/\rM$
(rotational).  The corresponding Jeans mass at that floor scales as
%
\begin{equation}
  M_J \propto T_{\rm floor}^{3/2} \rho^{-1/2}
        \propto (\ra \rmm)^3 \rM^{-3/2}\,\rho^{-1/2},
\end{equation}
%
so a heavier dark proton, weaker coupling, or lighter dark electron
all yield a lower cooling floor and \emph{smaller} first-object masses
--- opening a window for sub-solar-mass dark black holes.  The
Chandrasekhar mass $M_{\rm Ch} \propto 1/M^2$ likewise shrinks with a
heavier dark proton, so LIGO's sub-solar-mass search directly probes
this parameter space \citep{Shandera2018}.

%----------------------------------------------------------------------
\section{Codes: modified Recfast++ and DarkKROME}
\label{sec:codes}

\paragraph{RecfastJulia (Paper~II).}
Public reimplementation of \texttt{Recfast++} in Julia, with all rate
coefficients rescaled and $\HDtwo$-network extensions.  Produces the
cosmological ionisation history $x_{\eD}(z)$, $x_{\HD}(z)$,
$x_{\HDtwo}(z)$ for user-specified $(M, m, \alphaD, \xi)$.  Output
serves as the initial condition for the collapse simulations of
Paper~III.

Location:
\href{https://github.com/jamesgurian/RecfastJulia}{github.com/jamesgurian/RecfastJulia}.

\paragraph{DarkKROME (Paper~III).}
Extension of the \texttt{KROME} chemistry solver \citep{Grassi2014}
that adds dark-sector species (prefix \texttt{Q}) and rescaled thermal
processes:
%
\begin{itemize}
  \item \texttt{-cooling=DARKATOM} --- atomic cooling
        (Compton, bremsstrahlung, collisional ionisation, excitation,
        recombination) with rates rescaled from the SM
        \texttt{ATOMIC}+\texttt{COMPTON}+\texttt{FF} pipeline.
  \item \texttt{-cooling=ADARKATOM} --- analytic (fully Rosenberg--Fan)
        atomic cooling, faster but slightly less accurate.
  \item \texttt{-cooling=DARKMOL} --- rovibrational $\HDtwo$ cooling.
  \item \texttt{-heating=DARKATOM}, \texttt{ADARKATOM},
        \texttt{DARKMOL} --- corresponding heating channels.
\end{itemize}
%
User specifies $\{M, m, \alphaD, \xi\}$ via \texttt{qp\_mass},
\texttt{qe\_mass}, \texttt{Dalpha}, \texttt{xi}.  Ships with a minimal
network \texttt{react\_dark} and a Python preprocessor that generates
the reaction subroutines from a text-file description.

The included one-zone spherical collapse test in Section~3 of
Paper~III reproduces (with SM parameters) the built-in
\texttt{CollapseZ} benchmark of KROME, and (with a shifted $\alphaD$
or $\xi$) matches the \citet{DAmico2018} ``mirror'' dark-matter
collapse phase diagram.

Location:
\href{https://bitbucket.org/mtryan83/darkkrome}{bitbucket.org/mtryan83/darkkrome}.

\paragraph{What's still missing.}
DarkKROME as of the paper release does not include the
$\HDtwo + \HDtwo \to 2\HD + \HDtwo$ heating channel that dominates at
$n \gtrsim 10^8\,{\rm cm^{-3}}$ in SM chemistry \citep{GalliPalla1998},
nor the three-body $\HDtwo$ formation reactions involving $\HD_{D,3}$
excited states.  Both channels are expected to matter deep into
runaway collapse ($n \gtrsim 10^8\,{\rm cm^{-3}}$) and are left for
future work.

%----------------------------------------------------------------------
\section{Full 3D cosmological simulations}
\label{sec:sims}

The one-zone DarkKROME calculations of Paper III test the chemistry in
a homogeneous collapsing sphere but do not evolve the dark fluid on a
mesh or resolve where in the cosmological box collapse actually
happens.  The next step is a full cosmological simulation.  This
section describes what such a simulation looks like structurally,
which parts differ from a standard $N$-body run, and the current
state-of-the-art (first zoom-in published in 2024).

\subsection{Which regime is which: collisionless vs.\ chemical}
\label{sec:regime_map}

Whether the dark component behaves as a collisionless species (like
CDM) or as a dissipative fluid (like SM baryons) depends on local
density.  Define the scattering count per particle per Hubble time,
%
\begin{equation}
  \mathcal{N}_{\rm scat}
    \sim n_D\,\sigma_{\rm eff}(v)\,v\,t_H,
\end{equation}
%
with $n_D$ the local dark-species density, $\sigma_{\rm eff}(v)$ an
effective (Coulomb + atomic + molecular) cross section, and $t_H$
Hubble time.  Two limits:
%
\begin{itemize}
  \item $\mathcal{N}_{\rm scat}\ll 1$: cosmological voids, outer
        haloes, low-density regions ($\rho_D \lesssim$ few
        $\rho_{\rm crit}$).  Dark particles rarely scatter; motion is
        driven by gravity alone.  Behaves like CDM for bulk dynamics.
        Cheapest solver: collisionless $N$-body.
  \item $\mathcal{N}_{\rm scat}\gg 1$: collapsed cores
        ($\rho_D \gtrsim 10^{4-8}\,\rho_{\rm crit}$).  Chemistry
        activates, atoms form, molecules form, gas cools.
        Behaves like primordial SM baryons.  Requires hydro + chemistry
        + cooling.
\end{itemize}
%
The transition happens over $\sim\!1$--$2$ decades in density.
Practical consequence: over $\gtrsim 99\%$ of a cosmological box the
dark fluid is effectively collisionless; only in $\lesssim 1\%$ of
the volume (dense cores) is the chemistry actively relevant.

\subsection{Three-component sim structure}
\label{sec:three_component}

A maximal atomic-DM cosmological simulation contains three
gravitationally-coupled species, each evolved with its own physics:

\begin{enumerate}
  \item \textbf{Standard-Model baryons.}  Full hydrodynamics with SM
        chemistry (H, He, H$_2$, dust, metals as appropriate),
        SM cooling curve $\Lambda_{\rm SM}(T,x_e,n)$, and star
        formation / stellar feedback in whichever subgrid model is
        adopted (FIRE-2, Illustris-TNG, etc.).
  \item \textbf{Collisionless CDM.}  A fraction $(1-\epsilon)$ of the
        total dark-matter mass is standard cold dark matter,
        represented as $N$-body particles evolved under gravity only.
  \item \textbf{Atomic-DM fluid.}  The dissipative fraction
        $\epsilon = \Omega_D/\Omega_{\rm DM}$ (typically $\sim 0.05$--
        $0.2$ from CMB $N_{\rm eff}$ and structure bounds) evolves as
        a hydrodynamic fluid with the DarkKROME chemistry network and
        rescaled cooling function $\Lambda_D(T_D,x_{\eD},n_D)$.
\end{enumerate}
%
The three species interact only via a shared gravitational potential:
%
\begin{equation}
  \nabla^2\Phi = 4\pi G\bigl(
      \rho_{\rm baryon}
    + \rho_{\rm CDM}
    + \rho_{\rm dark-fluid}
    \bigr),
  \label{eq:combined_poisson}
\end{equation}
%
with SM electromagnetism confined to species 1 and dark $U(1)$
confined to species 3 (assuming no kinetic-mixing coupling, which is
observationally constrained to $\lesssim 10^{-8}$).  No direct
energy exchange between SM baryons and dark atoms.  The two chemistry
networks are completely disjoint.

\paragraph{DM-only sub-case.}  Drop the SM-baryons species: two-fluid
sim with CDM particles + atomic-DM hydro.  Cheaper and appropriate
for studies focused on dark halo internal structure rather than
galaxy formation.

\subsection{Code components needed}
\label{sec:code_components}

Building an atomic-DM cosmological code requires assembling the
following pieces, most of which already exist in state-of-the-art
astrophysical hydrodynamics codes (GIZMO, AREPO, RAMSES, GADGET) and
just need atomic-DM-specific customisation:

\begin{itemize}
  \item \textbf{Hydro solver.}  SPH, moving mesh, meshless finite mass,
        or grid-based AMR.  Same solver can handle both SM baryons
        (species 1) and dark fluid (species 3) — just labelled
        differently.  The physics is the same Euler equations + a
        different cooling function per species.
  \item \textbf{Collisionless $N$-body integrator.}  Tree, FMM, or
        particle-mesh for CDM (species 2).
  \item \textbf{Shared gravity solver.}  Sums the mass from all three
        species into $\Phi$ via
        Eq.~\eqref{eq:combined_poisson}; each species pulls $-\nabla\Phi$
        as its gravitational acceleration.
  \item \textbf{Chemistry networks.}  One for SM primordial + metal
        chemistry (KROME, Grackle) and one for dark chemistry
        (DarkKROME).  Separate runtime state per species.
  \item \textbf{Cooling functions.}  $\Lambda_{\rm SM}$ from standard
        libraries; $\Lambda_D$ from the rescaled tables of
        \citet{Ryan2022a}.  Each cooling is a source term in its own
        species' internal-energy equation.
  \item \textbf{Cosmological initial conditions.}  Modified MUSIC2 or
        monofonIC to generate the two-fluid ($\delta_{\rm CDM}$,
        $\delta_{\rm dark-atom}$) initial displacement/velocity field
        with the DAO-suppressed dark transfer function
        \citep{Gurian2022}.  Initial ionisation fractions
        $x_{\eD}(z_{\rm start})$, $x_{\HD}(z_{\rm start})$,
        $x_{\HDtwo}(z_{\rm start})$ come from the cosmological
        recombination code (\texttt{RecfastJulia}).
\end{itemize}

\subsection{Existing cosmological zooms (Roy+2023, +2024; Gemmell+2024)}
\label{sec:roy2024}

Three published cosmological zoom-in simulations to date implement
the three-fluid architecture (SM baryons + CDM particles + dissipative
aDM fluid).  All use \textsc{gizmo} meshless-finite-mass hydrodynamics
with FIRE-2 galaxy-formation subgrid physics for the baryons.

\paragraph{Roy et al.\ 2023 (\citealt{Roy2023}).} First zoom of a
Milky-Way-mass ($M_{\rm halo}\sim 10^{12}\Msun$) galaxy with a
strongly-dissipative aDM sub-component.  Demonstrated that the dark
gas cools into a rotationally-supported \emph{dark disk} at the halo
centre, and that the disk fragments into ``dark clumps'' ---
sub-resolution clusters of dark compact objects.  Established the
qualitative behaviour of the pipeline.

\paragraph{Gemmell et al.\ 2024 (\citealt{Gemmell2023}).} Zooms of MW-analogue
hosts with focus on the \emph{substructure} response: satellite
galaxies inside aDM-modified MW-mass hosts have different mass
functions, concentrations, and spatial distributions than in CDM.
Shows aDM back-reacts on baryonic satellite properties, not just on
DM ones.

\paragraph{Roy et al.\ 2024 (\citealt{Roy2024}, ``Aggressively-Dissipative
Dark Dwarfs'').} First zoom targeting \emph{isolated dwarf galaxies}
($M_{\rm halo}\sim 10^{10}\Msun$).  Suite of seven simulations sweeping
aDM microphysics.  Setup:
%
\begin{itemize}
  \item \textbf{Code:} GIZMO meshless-finite-mass with FIRE-2 galaxy
        formation physics (stellar feedback, radiative transfer, ISM
        chemistry) for the baryons; a second hydro species for aDM
        with its own dark cooling.
  \item \textbf{Target:} FIRE-2 dwarf zoom-in suite at
        $M_{\rm halo}(z=0)\sim 10^{10}\Msun$.
  \item \textbf{Dark parameters:} atomic-DM fraction
        $f' = \Omega_{\rm aDM}/\Omega_{\rm DM} = 0.06$, dark-to-SM
        temperature ratio $\xi' = T_{\gD}'/T_{\rm CMB} < 0.5$ (CMB
        bound from \citealt{Bansal2022}).  Seven variants of
        $(m_{e'}, m_{p'}, \alpha')$ scanning ``fast'' vs ``slow''
        cooling regimes.
  \item \textbf{Dark cooling function (equilibrium):} adopts the
        Rosenberg-Fan 2017 rates for atomic processes (dark
        ionisation, excitation, recombination, bremsstrahlung)
        and includes 2-body molecular scattering channels
        $\HDtwo\text{-}\eD$, $\HDtwo\text{-}\pD$,
        $\HDtwo\text{-}\HD$, $\HDtwo\text{-}\HDtwo$ from
        \citet{Ryan2022a,Ryan2022b}.  Species number densities are
        computed under \emph{collisional-ionisation equilibrium}
        (ionisation and recombination assumed to equilibrate faster
        than dynamical time), so no non-equilibrium chemistry
        network is evolved on-the-fly.  This is analogous to using
        an equilibrium cooling table
        $\Lambda(T, n, Z)$ for SM baryons instead of running
        KROME/Grackle non-equilibrium rate equations.
\end{itemize}
%
Findings:
%
\begin{itemize}
  \item In the ``aggressively cooling'' corner of parameter space,
        aDM gas cools rapidly at high redshift, most of the aDM in
        the halo collapses centrally into a rotationally-supported
        dark disk, and the disk fragments into ``dark clumps''
        (unresolved clusters of dark compact objects).
  \item Inner halo relaxes to an approximately \emph{isothermal}
        distribution described by a simple fitting function ---
        universal across the parameter range that produces aggressive
        cooling.
  \item Even at $f'\sim 5$--$6\%$, central densities of classical
        dwarfs are enhanced by \emph{more than an order of magnitude}
        relative to CDM+FIRE-2 baseline runs.  A sharp observational
        target.
  \item Baryonic disc shrinks (deeper potential from dark disk) and
        stellar age distribution skews older.  Even a subdominant
        aDM component has O(1) effects on the visible galaxy.
\end{itemize}
%
The natural extension is to full-box (not zoom) cosmological
simulations at similar resolution, needed for statistical
predictions rather than case studies.  Also open: extending the
molecular network to include three-body $\HDtwo$ formation and the
$\HDtwo + \HDtwo$ heating channel that dominates at
$n\gtrsim 10^8\,{\rm cm^{-3}}$ --- currently absent from both
DarkKROME (Paper III) and the Roy+2024 network, so runaway collapse
into a dark protostar remains resolution-limited.

\subsection{What's still open}
\label{sec:open}

\begin{itemize}
  \item Cosmological box (not zoom) sims with all three components ---
        needed for statistical predictions (halo mass function of
        dissipative-DM haloes, dark-star luminosity function, dark-BH
        merger rate distribution).
  \item Extension of DarkKROME to include $\HDtwo + \HDtwo$ heating and
        $\HD_{D,3}$ excited-state reactions at
        $n\gtrsim 10^8\,{\rm cm^{-3}}$, needed to follow
        runaway collapse into a dark protostar or dark black hole.
  \item Radiative feedback from dark protostars: do they emit a
        Str\"omgren volume of dark photons that ionise nearby
        dark gas, or do they accrete quietly?
  \item Joint dark + baryonic feedback: does dark-DM cooling
        change SM star formation efficiencies enough to alter dwarf
        stellar-mass functions?  Roy+2024 suggests yes, but the
        interaction has not been mapped systematically.
\end{itemize}

%----------------------------------------------------------------------
\section{Observational hooks}
\label{sec:observational}

\begin{itemize}
  \item \textbf{Linear-scale $P(k)$ suppression.}  Dark diffusion +
        DAO produce a cutoff in the transfer function analogous to WDM
        with an oscillatory feature; galaxy surveys (BOSS, DESI, Euclid),
        Lyman-$\alpha$ forest, and 21-cm at cosmic dawn all constrain
        the cutoff scale.
  \item \textbf{Sub-solar mass compact objects.}  If dark hydrogen
        cooling delivers Jeans masses down to $\sim 10^{-3}\,M_\odot$
        for some $(M, m)$, the resulting dark black holes are
        sub-Chandrasekhar for the SM.  LIGO/Virgo/KAGRA constrain such
        binaries via gravitational-wave inspirals with characteristic
        masses forbidden to standard-model stellar remnants.
  \item \textbf{Halo internal structure.}  Dissipative dark matter
        can form thin discs and rotationally supported subhalos
        \citep{Fan2013}, altering satellite kinematics and stream
        heating budgets --- relevant to small-scale CDM problems.
\end{itemize}

%----------------------------------------------------------------------
\section{Summary and open questions}
\label{sec:summary}

The Ryan/Gurian/Shandera/Jeong trilogy is best read as a self-contained
pipeline:
\begin{enumerate}
  \item Paper~I: rescale every atomic/molecular rate coefficient from
        SM hydrogen to dark hydrogen, based on Born--Oppenheimer
        structure.
  \item Paper~II: use those rates to run cosmological
        recombination + molecule formation, and evaluate the
        diffusion/DAO cutoff of the linear transfer function.
  \item Paper~III: feed the cosmological freeze-out abundances into
        DarkKROME one-zone collapse simulations of individual dark
        halos, tracking cooling and fragmentation as a function of
        $(M, m, \alphaD, \xi)$.
\end{enumerate}

Open directions signposted by the papers:
\begin{itemize}
  \item Full $N$-body + hydrodynamic simulations with dark chemistry
        --- currently only one-zone tests exist.
  \item Extension of the network to $\HD_{D,3}$ excited states and
        collisional dissociation at $n > 10^8\,{\rm cm^{-3}}$, needed
        to track the last decades of runaway collapse.
  \item Baryonic feedback in dark protostars --- do they radiate,
        ionise a Str\"omgren volume, or accrete quietly?
  \item Statistical prior on $(M, m, \alphaD, \xi)$ from a physical
        model of the dark sector origin.
\end{itemize}

The overall lesson is that dissipative dark matter is a
\emph{calculable} extension of CDM: the chemistry that governs Pop III
star formation in the SM has a clean parametric analogue in a dark
sector with atoms, and the resulting compact-object mass spectrum is a
direct observational target for current gravitational-wave detectors.

%======================================================================
\appendix

\section{Hydrogenic atom: where $\alpha$, $a_0$, $E_H$ come from}
\label{app:hydrogenic}

The atomic-DM model of \S\ref{sec:model} imports the parameters
$(a_{0,D}, E_{H_D}, \alpha_D)$ from SM hydrogen by rescaling.  This
appendix derives those quantities from first principles for a
$U(1)$ gauge theory, so the ``dark'' versions are transparent.

\subsection{Fine-structure constant from $U(1)$ gauge coupling}
\label{app:alpha}

Consider two oppositely charged fermions of charges $\pm e$
(``electron'' and ``proton'') interacting via a massless $U(1)$
gauge field.  The Lagrangian contains
%
\begin{equation}
  \mathcal{L}_{\rm gauge}
    = -\tfrac{1}{4}F_{\mu\nu}F^{\mu\nu}
      + \bar\psi(i\gamma^\mu D_\mu - m)\psi,
  \qquad
  D_\mu = \partial_\mu - ie A_\mu,
\end{equation}
%
with fermion covariant derivative $D_\mu$ and coupling constant $e$.
Instantaneously (Coulomb-gauge, non-relativistic limit) the potential
between two charges separated by $r$ is
%
\begin{equation}
  V(r) = \frac{e^2}{4\pi\epsilon_0}\,\frac{1}{r}
       \equiv \frac{\alpha \hbar c}{r}
  \qquad\text{(SI units)},
\end{equation}
%
which \emph{defines} the fine-structure constant
%
\begin{equation}
  \alpha \equiv \frac{e^2}{4\pi\epsilon_0\,\hbar c}.
\end{equation}
%
Numerically $\alpha \simeq 1/137$ for the SM.  In natural units
$\hbar = c = \epsilon_0 = 1$, this simplifies to
$\alpha = e^2/(4\pi)$ --- $\alpha$ is just the squared gauge coupling
in the Feynman-rule normalisation of QED.

For a dark $U(1)$ with coupling $g_D$ (i.e.\ dark charges
$\pm g_D$), the same construction defines
%
\begin{equation}
  \alpha_D \equiv \frac{g_D^2}{4\pi}
  \qquad\text{(natural units)}.
\end{equation}
%
$\alpha_D$ is a free parameter of the model; SM hydrogen corresponds
to $\alpha_D = \alpha_{\rm EM} = 1/137$.  The rescaling ratio
$r_\alpha \equiv \alpha_D/\alpha_{\rm EM}$ used throughout the trilogy
is the ratio of dark-to-SM gauge coupling squared.

\subsection{Bohr radius from Schr\"odinger equation}
\label{app:a0}

For a hydrogen-like bound state $(\pm e)$ around $(\mp e)$ with
reduced mass $\mu \simeq m$ (electron mass, since proton $\gg$
electron), the non-relativistic Schr\"odinger equation is
%
\begin{equation}
  \left[-\frac{\hbar^2}{2m}\nabla^2 - \frac{\alpha\hbar c}{r}\right]\psi
    = E\psi.
  \label{eq:hydrogen_SE}
\end{equation}
%
Dimensional analysis: the only length that can be built from
$(m,\hbar,\alpha c)$ is
%
\begin{equation}
  a_0 = \frac{\hbar}{m c\,\alpha}
      = \frac{\hbar^2}{m\,(\alpha\hbar c)}
      = 5.29\times 10^{-9}\,{\rm cm}
      \qquad\text{(SM)}.
  \label{eq:bohr_radius}
\end{equation}
%
This is the Bohr radius; solving \eqref{eq:hydrogen_SE} exactly
(hypergeometric functions, ground-state $\psi_{100}\propto e^{-r/a_0}$)
confirms this scale.  In natural units $a_0 = 1/(m\alpha)$.

For dark hydrogen with electron mass $m$ and coupling $\alpha_D$,
%
\begin{equation}
  a_{0,D} = \frac{1}{m\,\alpha_D}
          = a_0\,\frac{1}{r_\alpha\,r_m},
\end{equation}
%
where $r_m \equiv m/m_e$.  Heavier dark electrons and stronger dark
coupling both shrink the atom.

\subsection{Rydberg (binding energy) from virial theorem}
\label{app:rydberg}

For the Coulomb potential $V(r) = -\alpha\hbar c/r$, the virial
theorem gives $\langle T\rangle = -\tfrac{1}{2}\langle V\rangle$, so
$E = T + V = \tfrac{1}{2}\langle V\rangle < 0$.  Evaluating $\langle
V\rangle$ on the ground-state $\psi_{100}$ requires
$\langle 1/r\rangle = 1/a_0$, giving
%
\begin{equation}
  E_1 = -\tfrac{1}{2}\,\alpha\hbar c\,\langle 1/r\rangle
      = -\frac{1}{2}\,\frac{\alpha\hbar c}{a_0}
      = -\tfrac{1}{2}\,m c^2 \alpha^2
      \simeq -13.6\eV.
  \label{eq:rydberg}
\end{equation}
%
$E_H \equiv |E_1| = \tfrac{1}{2}m c^2\alpha^2$ is the Rydberg energy;
in natural units $E_H = m\alpha^2/2$.  For dark hydrogen:
%
\begin{equation}
  E_{H_D} = \tfrac{1}{2}\,m\,\alpha_D^2
    = E_H\,r_\alpha^2\,r_m.
\end{equation}
%
\emph{Both $a_0$ and $E_H$ depend only on the fermion mass $m$ and
the coupling $\alpha$ --- not on the ``proton'' mass $M$.}  The
proton enters only via (i) the reduced mass correction
$\mu = mM/(m+M)\simeq m$ when $M\gg m$, and (ii) hyperfine splitting
$\propto m/M$ which we ignore.  This is why the Ryan+2022 rescaling
is one-parameter in $r_\alpha$ and $r_m$: molecular energy scales are
electronic, so the proton mass only enters molecular \emph{rotation}
and \emph{vibration}, which are ``adiabatically slow'' compared to
electronic transitions.

\subsection{Fine structure and orders of magnitude}
\label{app:orders}

The atomic-DM chemistry pipeline treats three energy scales:
%
\begin{equation}
  E_{\rm rot} \ll E_{\rm vib} \ll E_H,
\end{equation}
%
scaling as
%
\begin{align}
  E_{\rm rot} &\sim \frac{\hbar^2}{M X_0^2}
      \sim \frac{m^2\alpha^2}{M}
      = E_H\,\frac{m}{M},
    &&\text{(rotational)}, \\
  E_{\rm vib} &\sim \hbar\sqrt{\frac{k}{M}}
      \sim \alpha^2\sqrt{\frac{m^3}{M}}
      = E_H\,\sqrt{\frac{m}{M}},
    &&\text{(vibrational)}, \\
  E_H  &= \tfrac{1}{2}m\alpha^2,
    &&\text{(electronic)}.
\end{align}
%
So molecular spectra separate cleanly into three well-separated
energy scales, ordered by powers of $\sqrt{m/M}$.  For SM hydrogen,
$m/M = 5.4\times 10^{-4}$, so $E_{\rm rot} \sim 10^{-3} E_H\sim
10^{-2}\eV$ (matches SM $H_2$ rotational lines at
$\sim 10^{-2}\eV$), $E_{\rm vib}\sim 10^{-2}E_H\sim 0.1\eV$
(matches SM $H_2$ vibrational spacings).  For heavier dark ``protons''
($r_M\gg 1$) or lighter dark ``electrons'' ($r_m\ll 1$), these scales
run in known ways --- see Table~1 of \citet{Ryan2022a} for the
complete rescaling table.

The Born--Oppenheimer approximation (nuclear motion adiabatic to
electronic motion) is what makes this three-scale hierarchy exact
in the $m/M \to 0$ limit, and is why Paper~I's rescaling is
theoretically rigorous rather than a fitting exercise.

%======================================================================
\section{Cosmology primer for ISM/star-formation readers}
\label{app:cosmo_primer}

Readers whose home ground is ISM chemistry, cooling functions, and
star formation may not carry the standard cosmological toolbox.  This
appendix collects the pieces that appear in the main text.  A fuller
treatment is in \texttt{cosmo\_ic.tex} in this repository, or Dodelson
\& Schmidt (2020), Baumann (2022), or Weinberg (2008).

\subsection{FRW background and the Hubble rate}

Cosmological spacetime is described by the Friedmann--Robertson--Walker
(FRW) metric, which for a spatially flat universe reads
%
\begin{equation}
  ds^2 = -c^2 dt^2 + a^2(t)\bigl[dr^2 + r^2 d\Omega^2\bigr].
\end{equation}
%
$a(t)$ is the \emph{scale factor}, normalised to $a(t_0) = 1$ today.
Physical distances between comoving points scale as $a(t)$: e.g.\ if
two galaxies sit at comoving separation $r$, their physical separation
today is $a(t_0)r = r$, but at earlier time $t$ was $a(t)r < r$.

The \emph{Hubble rate} is
%
\begin{equation}
  H(t) \equiv \frac{\dot a}{a},
\end{equation}
%
today $H_0 \approx 70\,{\rm km\,s^{-1}\,Mpc^{-1}}$, i.e.\ $H_0^{-1}
\approx 14\,{\rm Gyr}$ (the age of the universe).  The Friedmann
equation
%
\begin{equation}
  H^2 = \frac{8\pi G}{3}\bigl(\rho_r + \rho_m + \rho_\Lambda\bigr),
\end{equation}
%
partitions the total energy density into radiation ($\rho_r\propto
a^{-4}$), matter ($\rho_m\propto a^{-3}$), and dark energy ($\rho_\Lambda
= {\rm const}$).  Redshift is the standard time variable:
$1 + z = 1/a(t)$, so $z = 0$ today, $z \to \infty$ at the Big Bang.
Cosmological parameters (Planck 2018):
$\Omega_m = 0.315$, $\Omega_\Lambda = 0.685$, $h = 0.674$,
$\Omega_b h^2 = 0.0224$, $n_s = 0.965$, $\sigma_8 = 0.81$.

\subsection{Matter--radiation equality and the standard eras}

Radiation dominates early ($z\gtrsim 3400$), matter dominates the
middle era, dark energy dominates late ($z\lesssim 0.3$).  The key
epochs:
%
\begin{center}\begin{tabular}{lll}
\toprule
Epoch & $z$ & Physics \\
\midrule
BBN               & $\sim 10^{10}$ & primordial nucleosynthesis \\
Matter--radiation eq.\ ($z_{\rm eq}$) & $3400$ & structure growth begins \\
SM recombination & $1090$ & $e^-+p\to H$; CMB decouples \\
Reionisation      & $\sim 6$--$10$ & first stars ionise IGM \\
Present ($z=0$)   & $0$ & --- \\
\bottomrule
\end{tabular}\end{center}

Recombination (SM) is what \citet{Peebles1968} solved with his
effective 3-level model, adapted by \texttt{Recfast++} into a fast
code.  Paper~II ports this whole machinery to the dark sector, so a
few reminders on the SM version:

\paragraph{Saha equilibrium.}  In equilibrium, the ionisation
fraction $x_e = n_e/n_H$ satisfies
%
\begin{equation}
  \frac{x_e^2}{1-x_e} = \frac{1}{n_b}\left(\frac{m_e T}{2\pi}\right)^{3/2}\!
    e^{-E_H/T},
\end{equation}
%
which is the balance between recombination and photoionisation when
radiation and matter share a temperature.  Saha predicts $x_e$ drops
rapidly to zero at $T\ll E_H$; \emph{actually} $x_e$ freezes out at
$\sim 10^{-3}$ because recombination cannot keep up with expansion.

\paragraph{Peebles' 3-level.}  Recombination proceeds through the
excited $n=2$ state (case B recombination) --- direct $e^-+p\to H(1s)+\gamma$
would produce a Lyman-$\alpha$ photon that immediately re-ionises
another neutral atom, so it's ineffective.  The rate-limiting step is
the escape of $n=2$ to $n=1$ via forbidden two-photon decay
$2s\to 1s + 2\gamma$ or by cosmological redshifting Ly$\alpha$
out of the resonance line.  Freeze-out happens because the Universe
expands faster than recombination can catch equilibrium.

For dark hydrogen, all these ingredients rescale by $r_\alpha$,
$r_m$, $r_M$ per \S\ref{sec:rescaling}, and the equilibrium and
freeze-out redshifts shift accordingly.  Because the dark radiation
temperature is $\xi \cdot T_{\rm CMB}$, the actual dark-sector
recombination redshift depends on both the intrinsic dark binding
energy $E_{H_D}$ and $\xi$.

\subsection{Comoving vs physical, sub-horizon vs super-horizon}

Two coordinate systems constantly used:
%
\begin{itemize}
  \item \textbf{Physical} $(r_{\rm phys}, t)$: what a lab would
        measure.  Grows with $a(t)$.
  \item \textbf{Comoving} $(x, t)$ with $r_{\rm phys} = a(t)\,x$:
        removes the trivial expansion.  Two galaxies stay at
        fixed $x$ if only Hubble flow acts.
\end{itemize}
%
The \emph{comoving Hubble horizon} $1/(aH)$ is the causal-contact
scale.  A mode of comoving wavenumber $k$ is \emph{super-horizon} if
$k \lesssim aH$ (not causally connected across itself) or
\emph{sub-horizon} if $k \gtrsim aH$.  Structure formation is
dominated by sub-horizon dynamics; the transfer function suppression
in atomic DM (\S\ref{sec:dao}) is imprinted around horizon crossing
when $k = aH$.

\subsection{Linear power spectrum and the transfer function}

Cosmological perturbations start Gaussian, described by the
\emph{primordial power spectrum}
%
\begin{equation}
  P_{\rm prim}(k) = A_s\left(\frac{k}{k_*}\right)^{n_s - 1},
\end{equation}
%
essentially scale-invariant ($n_s\simeq 0.965$).  The
\emph{transfer function} $T(k)$ converts primordial to late-time:
%
\begin{equation}
  P_{\rm late}(k, z) = A_s\,T^2(k)\,D^2(z)\,\left(\frac{k}{k_*}\right)^{n_s-1},
\end{equation}
%
where $D(z)$ is the linear growth factor
(nearly $\propto a$ during matter domination).  $T(k)$ is what
distinguishes CDM from atomic DM (dark diffusion cutoff + DAOs), WDM
(free-streaming cutoff), or FDM (quantum-pressure cutoff at the Jeans
scale).

\subsection{$N_{\rm eff}$ and dark radiation bounds}

The Planck constraint on the effective number of relativistic
species,
%
\begin{equation}
  N_{\rm eff} = N_{\rm eff}^{\rm SM} + \Delta N_{\rm eff},
  \qquad
  \Delta N_{\rm eff} \lesssim 0.3\ (2\sigma),
\end{equation}
%
limits how much energy density the dark sector can carry in
relativistic species at CMB decoupling ($z\sim 1090$).  For a dark
sector at temperature $T_D = \xi T_\gamma$ with the standard
degrees-of-freedom counting, $\xi \lesssim 0.4$ is the safe bound.
This is why Paper~III adopts $\xi = 0.01$: to suppress dark-acoustic
oscillations well below observational detection while satisfying
$N_{\rm eff}$ trivially.

%======================================================================
\section{Particle-physics primer}
\label{app:particle_primer}

A few things worth clarifying if particle physics is not your
home ground.

\subsection{Natural units ($\hbar = c = 1$)}
\label{app:natural_units}

Particle physicists routinely ``set $\hbar = c = 1$'' --- meaning
Planck's constant and the speed of light are chosen as the units of
action and velocity, so all quantities are expressed in a single
mass (or equivalently energy) unit.  In these units:
%
\begin{itemize}
  \item Energy $=$ mass $=$ inverse length $=$ inverse time.
  \item $1\,{\rm GeV} = 5.07\times 10^{13}\,{\rm cm^{-1}}
        = 1.52\times 10^{24}\,{\rm s^{-1}}$.
  \item Bohr radius $a_0 = \hbar/(mc\alpha) \to 1/(m\alpha)$ (mass
        with inverse-length meaning).
  \item Rydberg energy $E_H = \tfrac{1}{2}mc^2\alpha^2\to
        \tfrac{1}{2}m\alpha^2$.
\end{itemize}
%
To recover conventional (SI or cgs) units, insert the missing
$\hbar$ and $c$ by dimensional analysis: e.g.\ $a_0 = 1/(m\alpha)$
becomes $\hbar/(mc\alpha)$ with dimensions restored.

\subsection{Why the atomic-DM model has exactly four parameters}
\label{app:four_params}

The model builds a dark hydrogenic $U(1)$ gauge theory (see
\S\ref{app:U1}).  The Lagrangian
$\mathcal L = -\tfrac{1}{4}F_{\mu\nu}^D F_D^{\mu\nu} + \bar\pD(i\gamma^\mu
D_\mu - M)\pD + \bar\eD(i\gamma^\mu D_\mu - m)\eD$
contains three fundamental constants of nature: the two fermion
masses ($M$, $m$) and the gauge coupling $g_D$ (equivalently
$\alpha_D = g_D^2/4\pi$).  No other parameter can appear: the
$U(1)$ dictates the interaction form, gauge invariance forbids
fermion-photon Yukawa couplings and mediator masses, and imposing
strict charge-neutrality kills the ``how much dark antimatter''
degree of freedom.  So the \emph{microphysics} is 3-dimensional:
$(M, m, \alpha_D)$.

The \emph{cosmological} setup adds one more parameter: at some early
time the dark sector is populated with radiation and matter, and its
radiation temperature $T_{\gD}$ may differ from the CMB temperature
$T_\gamma$.  This is captured by
%
\begin{equation}
  \xi \equiv T_{\gD}/T_\gamma.
\end{equation}
%
$\xi$ is an initial-condition parameter --- it depends on how the
dark sector was populated during reheating --- rather than a
Lagrangian parameter.

Total: $(M, m, \alpha_D, \xi)$ specifies the full model.  The
overall dark-matter abundance $\Omega_D$ is fixed by requiring
observed structure; it is not an independent knob.

\subsection{Cross-section units: astro vs particle physics}

Astrophysicists measure ``interaction strength'' via $\sigma/m$ in
$\cmg$; particle physicists use $\mathrm{cm}^2/\mathrm{GeV}$ or barns.
Conversion:
%
\begin{equation}
  1\cmg = 1.78\times 10^{-24}\,\mathrm{cm}^2\,\mathrm{GeV}^{-1}
       = 1.78\,\mathrm{barn}\,\mathrm{GeV}^{-1}.
\end{equation}
%
For reference:
\begin{itemize}
  \item Weak cross section (WIMP): $\sim 10^{-38}\cmg$
  \item Electromagnetic (Thomson): $\sim 10^{-25}\cmg$
  \item Strong (n-p): $\sim 10^{-1}\cmg$
  \item Astrophysical SIDM target: $\sim 1\cmg$
\end{itemize}

\subsection{Gauge symmetry and $U(1)$}
\label{app:U1}

``$U(1)$'' is the group of unit complex numbers under multiplication:
$U(1) = \{e^{i\theta}\,:\,\theta\in\mathbb R\}$ ($U$ for
\emph{unitary}, $1$ for one-dimensional).  Physically, it's the
symmetry group of a phase rotation: rotate every quantum field by
$\psi\to e^{i\theta}\psi$ globally, and observable physics doesn't
change.

The step from \emph{global} to \emph{local} $U(1)$ gauge symmetry
(where $\theta = \theta(\mathbf{x},t)$ varies point-by-point) is
Weyl's / Yang--Mills' insight: to preserve invariance, the
derivative $\partial_\mu\psi$ must be replaced by a
\emph{covariant derivative}
$D_\mu\psi = (\partial_\mu - ie A_\mu)\psi$ that ``absorbs'' the
spatially varying phase.  This forces the introduction of a gauge
field $A_\mu$ (the photon in QED, or a dark photon in the atomic-DM
model) that transforms as $A_\mu\to A_\mu + \tfrac{1}{e}\partial_\mu\theta$
under the symmetry.  In other words:
\emph{local $U(1)$ symmetry $\Rightarrow$ existence of a massless
gauge boson}, plus a specific structure of interaction between the
boson and its charged particles.  Cross-fertilisation of this idea
with $SU(2)\times U(1)$ (weak+electromagnetic) and $SU(3)$
(strong) gives the Standard Model.

For atomic dark matter we take the simplest possible dark-sector
extension: a fresh $U(1)$, its own gauge boson $A_\mu^D$ (the
``dark photon''), its own coupling $g_D$ (a free parameter),
and two fermions of opposite dark charge $\pm g_D$ (the ``dark
proton'' and ``dark electron'').  The SM particles are dark-neutral,
and the two sectors decouple electromagnetically.  Everything else
--- atoms, molecules, chemistry --- reappears in the dark sector by
the same dynamics that produced it in the SM, with $\alpha \to
\alpha_D$ and $(m_e, m_p) \to (m, M)$.  This is why the whole
Ryan+2022 rescaling programme works: it's the same physics, one gauge
theory reused.

\subsection{Dark $U(1)$ vs QED}

Standard-model quantum electrodynamics (QED) has coupling
$\alpha_{\rm EM} \simeq 1/137$; every charged particle interacts with
the photon via $e A_\mu \bar\psi\gamma^\mu\psi$.  A ``dark $U(1)$'' is
the same mathematical structure with a separate gauge field
$A_\mu^D$, its own coupling $g_D$, and its own charged fermions
($\pD, \eD$).  The Standard Model particles are (by construction)
neutral under this new $U(1)$, so dark and visible sectors do not
interact electromagnetically --- except through kinetic mixing
$\epsilon\,F_{\mu\nu}F_D^{\mu\nu}$, which is small ($\epsilon\lesssim
10^{-8}$) enough to be irrelevant for our cosmology but is the target
of direct-detection experiments.

\subsection{Born approximation}

Cross sections at low velocity $v$ are typically computed using the
\emph{Born approximation}: treat the interaction as a small
perturbation on a plane-wave incoming state.  For an elastic Yukawa
potential $V(r) = \alpha e^{-m_\phi r}/r$, the Born differential
cross section is
%
\begin{equation}
  \frac{d\sigma}{d\Omega}
    = \frac{4\alpha^2 m^2}{[q^2 + m_\phi^2]^2},
  \qquad
  q = 2m v\sin(\theta/2),
\end{equation}
%
where $\vec q$ is the momentum transfer.  Two limits:
%
\begin{itemize}
  \item $q\gg m_\phi$: $d\sigma/d\Omega \propto 1/q^4 \propto 1/v^4$
        (Rutherford-like, high velocity).
  \item $q\ll m_\phi$: $d\sigma/d\Omega \propto 1/m_\phi^4$
        (velocity-independent, low velocity).
\end{itemize}
%
This is why dark-photon-mediated SIDM has strong velocity dependence:
the mediator mass sets the crossover scale.

\bibliography{references}

\end{document}
